\documentclass[12pt,letterpaper]{article}
\usepackage[utf8]{inputenc}
\usepackage[T1]{fontenc}
\usepackage[margin=2.54cm]{geometry}
\usepackage{mathptmx}
\usepackage[table]{xcolor}
\usepackage{array}
\usepackage[hidelinks]{hyperref}
\pagestyle{empty}
\color{black}
\renewcommand{\arraystretch}{1.12}
\definecolor{cab}{gray}{0.75}
\newcolumntype{A}{>{\raggedright\arraybackslash}p{3.15cm}}
\newcolumntype{B}{>{\raggedright\arraybackslash}p{7.65cm}}
\newcolumntype{C}{>{\raggedright\arraybackslash}p{4.35cm}}
\begin{document}
\begin{titlepage}
\centering
\vspace*{1cm}
{\Large \textbf{UNIVERSIDAD DE CÓRDOBA}}\\[0.3cm]
{\large FACULTAD DE INGENIERÍA}\\[0.2cm]
{\large PROGRAMA DE INGENIERÍA DE SISTEMAS}\\[0.2cm]
{\large SEDE MONTERÍA}\\[1.7cm]
\rule{\textwidth}{0.5pt}\\[0.3cm]
{\Huge \textbf{FICHA DE SOFTWARE}}\\[0.4cm]
{\LARGE \textbf{SolarBell}}\\[0.3cm]
\rule{\textwidth}{0.5pt}\\[0.5cm]
{\large \textit{Ficha técnica de las funcionalidades de software del sistema SolarBell.}}\\[2.5cm]
\begin{minipage}[t]{0.42\textwidth}
\begin{flushleft}
\textbf{Presentado por:}\\[0.2cm]
Andrés Felipe Ortiz Rodríguez,\\
Néstor Junior Hernández López,\\
José Enrique Otero Pineda,\\
Julián Camilo Doria Guevara
\end{flushleft}
\end{minipage}\hfill
\begin{minipage}[t]{0.52\textwidth}
\begin{flushright}
\textbf{Docentes:}\\[0.2cm]
Antonio Fernando Movilla Quintero\\
\textit{Sistemas Operativos}\\[0.3cm]
Milton Eduardo Hernández Zakzuk\\
\textit{Arquitectura del Computador}
\end{flushright}
\end{minipage}\\[2.5cm]
\vfill
Montería, Córdoba\\
Noviembre de 2026
\end{titlepage}
\newpage
\noindent{\large \textbf{Ficha de Software}}\\[0.25cm]
\noindent\begin{tabular}{|A|B|C|}
\hline
\rowcolor{cab}\multicolumn{2}{|l|}{\textbf{Nombre del Proyecto:}} & \textbf{Plataforma:} \\
\hline
\multicolumn{2}{|l|}{SolarBell: Sistema de timbre automatizado con energía solar} & Móvil (solo Android) \\
\hline
\rowcolor{cab}\textbf{Funcionalidad} & \multicolumn{2}{l|}{\textbf{Descripción}} \\
\hline
Conectividad & \multicolumn{2}{>{\raggedright\arraybackslash}p{11.55cm}|}{La aplicación habla con la campana por HTTP, en la red Wi-Fi SolarBell o en la red Wi-Fi del colegio, sin internet y sin servidor. Consulta su estado y le envía el horario y las órdenes.} \\
\hline
Puesta en Hora & \multicolumn{2}{>{\raggedright\arraybackslash}p{11.55cm}|}{Al conectarse, pone la campana en la hora y la zona horaria del teléfono. El reloj del dispositivo conserva esa hora si se corta la energía.} \\
\hline
Gestión de Horarios & \multicolumn{2}{>{\raggedright\arraybackslash}p{11.55cm}|}{Permite consultar y modificar el horario semanal: nombre, hora, cantidad de toques, duración y pausa de cada toque, y los días activos de lunes a viernes. El horario completo queda en la memoria de la campana y ella lo ejecuta sola. Si está fuera de alcance, el cambio no se guarda.} \\
\hline
Fechas Especiales & \multicolumn{2}{>{\raggedright\arraybackslash}p{11.55cm}|}{Permite definir una fecha que ese día reemplaza el horario semanal, con timbres propios o sin ninguno.} \\
\hline
Consulta del Día & \multicolumn{2}{>{\raggedright\arraybackslash}p{11.55cm}|}{Muestra la fecha de hoy, el tiempo que falta para el próximo timbre, los timbres del día y las fechas especiales próximas.} \\
\hline
Toque Manual & \multicolumn{2}{>{\raggedright\arraybackslash}p{11.55cm}|}{Permite hacer sonar la campana en el momento, eligiendo las veces, la duración y la pausa de cada toque, y detener el toque en curso. El tiempo lo lleva la campana.} \\
\hline
Estado de la Campana & \multicolumn{2}{>{\raggedright\arraybackslash}p{11.55cm}|}{Muestra si la campana está al alcance, si su reloj está en hora y la versión del programa que ejecuta.} \\
\hline
\end{tabular}
\end{document}
